\documentclass[11pt,a4paper]{article}
\usepackage[T1]{fontenc}
\usepackage[utf8]{inputenc}
\usepackage[main=italian,provide=*]{babel}
\usepackage{amsmath,amssymb}
\usepackage{geometry}
\geometry{margin=2.3cm,top=2.5cm,bottom=2.7cm}
\usepackage{booktabs}
\usepackage[table]{xcolor}
\usepackage{tcolorbox}
\tcbuselibrary{skins,breakable}
\usepackage{titlesec}
\usepackage{enumitem}
\usepackage{seqsplit}
\usepackage{needspace}
\usepackage{fancyhdr}
\usepackage{hyperref}
\hypersetup{colorlinks=true,linkcolor=Sepia,citecolor=Sepia,urlcolor=Sepia}

\definecolor{inkblue}{HTML}{1B3A5C}
\definecolor{lightblue}{HTML}{EAF1F8}
\definecolor{accent}{HTML}{B0562A}
\definecolor{softgray}{HTML}{6B6B6B}
\definecolor{okgreen}{HTML}{2E6B3E}
\definecolor{okgreenbg}{HTML}{EAF5EC}
\definecolor{warnamber}{HTML}{9C6B12}
\definecolor{warnbg}{HTML}{FBF1E0}
\definecolor{advisorpurple}{HTML}{5B3A7A}
\definecolor{advisorbg}{HTML}{F1ECF6}
\definecolor{notebar}{HTML}{9C6B12}

\titleformat{\section}
  {\normalfont\Large\bfseries\color{inkblue}}
  {\thesection}{0.6em}{}[{\color{inkblue!40}\titlerule}]
\titlespacing{\section}{0pt}{0.8em}{0.4em}
\titleformat{\subsection}
  {\normalfont\large\bfseries\color{accent}}
  {}{0em}{}
\titlespacing{\subsection}{0pt}{0.5em}{0.2em}

\pagestyle{fancy}
\fancyhf{}
\renewcommand{\headrulewidth}{0.4pt}
\fancyhead[R]{\small\color{softgray}\thepage}
\fancyfoot[C]{\small\color{softgray}Tesi triennale, Samuele Schiavi --- Relatore: Prof.\ Alessandro Chiesa}

\newtcolorbox{keyeq}[1][]{
  colback=lightblue, colframe=inkblue!70, boxrule=0.6pt,
  arc=2pt, left=8pt, right=8pt, top=4pt, bottom=4pt, #1
}
\newtcolorbox{panelbox}[1]{
  colback=white, colframe=softgray!50, boxrule=0.5pt,
  arc=2pt, left=7pt, right=7pt, top=3pt, bottom=3pt,
  fonttitle=\bfseries\color{inkblue}, title={#1}
}
\newtcolorbox{okbox}[1]{
  colback=okgreenbg, colframe=okgreen!60, boxrule=0.6pt,
  arc=2pt, left=7pt, right=7pt, top=3pt, bottom=3pt,
  fonttitle=\bfseries\color{okgreen}, title={#1}
}
\newtcolorbox{warnbox}[1]{
  colback=warnbg, colframe=warnamber!70, boxrule=0.6pt,
  arc=2pt, left=7pt, right=7pt, top=3pt, bottom=3pt,
  fonttitle=\bfseries\color{warnamber}, title={#1}
}
\newtcolorbox{advisorbox}[1]{
  colback=advisorbg, colframe=advisorpurple!60, boxrule=0.6pt,
  arc=2pt, left=7pt, right=7pt, top=4pt, bottom=4pt,
  fonttitle=\bfseries\color{advisorpurple}, title={#1}
}
\newtcolorbox{editnote}{
  colback=white, colframe=white, boxrule=0pt,
  borderline west={2.2pt}{0pt}{notebar},
  arc=0pt, left=10pt, right=4pt, top=2pt, bottom=2pt,
  fontupper=\itshape\small\color{softgray!90!black}
}
\fancyhead[L]{\small\color{softgray}Come si dividono i compiti (anello): analisi\_espressivita\_PMA vs confronto\_ansatz\_entangler}

\title{\vspace{-1.5em}\color{inkblue}Come si dividono i compiti (anello)\\[0.15em]
\Large\normalfont\color{softgray}\texttt{analisi\_espressivita\_PMA\_anello.ipynb} e \texttt{confronto\_ansatz\_entangler\_trimero\_anello.ipynb}}
\author{Note di lavoro --- Tesi triennale, Samuele Schiavi\\ Relatore: Prof.\ Alessandro Chiesa}
\date{}

\begin{document}
\sloppy
\maketitle
\thispagestyle{fancy}
\vspace{-2em}

\begin{editnote}
A differenza del dimero (dove l'indagine mirata precede e viene poi
importata dal confronto sistematico), qui l'ordine cronologico reale è
\textbf{invertito}: il residuo sotto DM emerge già nella lettura dei
risultati di \texttt{confronto\_ansatz\_entangler\_trimero\_anello.ipynb}
(riportato ma non spiegato), ed è quell'osservazione a motivare
\emph{a posteriori} l'indagine dedicata. La divisione dei compiti resta
la stessa nella sostanza --- una domanda di design isolata da un lato,
il quadro sistematico dall'altro --- ma non nella sequenza temporale.
\end{editnote}

\begin{center}
\renewcommand{\arraystretch}{1.5}
\small
\begin{tabular}{@{}p{2.0cm}p{5.7cm}p{5.7cm}@{}}
\toprule
\rowcolor{inkblue!10}
 & \textbf{\seqsplit{analisi\_espressivita\_PMA\_anello.ipynb}} & \textbf{\seqsplit{confronto\_ansatz\_entangler\_trimero\_anello.ipynb}}\\
\midrule
\textbf{Domanda} &
Il potere espressivo delle famiglie PMA regge davvero sotto DM, non
solo a $D=0$? (la domanda "dove mettere le rotazioni" è già risolta nel
dimero e qui semplicemente ereditata, non riaperta) &
Come si comporta l'intera famiglia di 10 ansatz su tutto l'asse $b$ e
al crescere di $D$?\\
\textbf{Metodo} &
Isolamento del punto peggiore dalla griglia standard ($b=0.05$, non
$b_c$ come atteso), poi verifica dedicata: 60 restart, ottimizzazione
diretta della fidelity (non dell'energia) su 5 varianti sospette &
Sweep sistematico su griglia $(b,D)$, 2 restart, confronto con 10
ansatz insieme (cache su file)\\
\textbf{Esito} &
La famiglia "2q" ha un \emph{tetto strutturale vero}
$\mathcal F\approx0.99831$, identico per 6 e 9 parametri (3 varianti
indipendenti convergono allo stesso valore); la famiglia "1q" --
inferiore a $D=0$ -- lo supera e raggiunge $\mathcal F=1$ esatto a 9
parametri &
\emph{Rileva} il sintomo (residuo $\approx0.995$--$0.998$ per quasi
tutte le famiglie PMA con DM) ma non lo diagnostica; individua
\texttt{D: Ry-RBS-Ry} come ansatz canonico per gate-count\\
\bottomrule
\end{tabular}
\end{center}

\begin{okbox}{In sintesi}
\texttt{confronto\_ansatz\_entangler\_trimero\_anello.ipynb} nasce prima
e mostra un numero senza spiegarlo --- la tabella riassuntiva riporta
il residuo sotto DM come un dato di fatto, non come un problema da
risolvere. \texttt{analisi\_espressivita\_PMA\_anello.ipynb} nasce
\emph{dopo}, proprio per rispondere a quella domanda rimasta aperta:
isola il punto peggiore, lo verifica con un'ottimizzazione molto più
approfondita (60 restart contro i 2 dello sweep standard), e scopre che
il residuo non è un artefatto ma un vero limite di espressività --
inatteso, perché va nella direzione opposta a quanto osservato a
$D=0$, dove la famiglia "2q" è sempre la più efficiente. Il notebook di
confronto è stato poi aggiornato con una nota che rimanda esplicitamente
a questa scoperta.
\end{okbox}

\end{document}
